% Originally: content/week-04.tex, \section{The Jacobi matrix} (Corsin Week 4, Monday)

% Source: Corsin Nick/Class Notes/Week 4.pdf, p. 3
\section{The Jacobi matrix}
\label{sec:jacobi_matrix}

Since $DF_{x_0} : \mathbb{R}^n \to \mathbb{R}^m$ (\cref{def:differentiable}) is a linear map, there is a \textbf{1--1
correspondence with a matrix once a basis is fixed}! In the standard basis of $\mathbb{R}^n$,
$(e_1,\dots,e_n)$, this is the \newterm{Jacobi matrix} \germanfor{Jacobi-Matrix}:
\[DF_x \;\cong\; \Jac F(x) = \begin{pmatrix} \dfrac{\partial F_1}{\partial x_1} & \cdots & \dfrac{\partial F_1}{\partial x_n} \\[1ex] \vdots & \ddots & \vdots \\[1ex] \dfrac{\partial F_m}{\partial x_1} & \cdots & \dfrac{\partial F_m}{\partial x_n}\end{pmatrix}\]
($m$ rows, $n$ columns), where we wrote $F(x) = (F_1, \dots, F_m)\transp$ and
\[\frac{\partial F}{\partial x_i} =: \partial_i F(x)\]
is the \newterm{$i$-th partial derivative} \germanfor{partielle Ableitung} at $x$:
\[\partial_i F(x) = \lim_{s\to 0}\frac{F(x+se_i)-F(x)}{s}.\]
This defines a function $\partial_i F : U \to \mathbb{R}^m$. If $F$ is differentiable, then
\[\partial_i F(x) = DF_x(e_i).\]

\begin{importantremark}
The partial derivatives can exist even if $F : U \to \mathbb{R}^m$ is not differentiable (recall \cref{ex:partial_derivatives_not_continuous})!
\end{importantremark}

% Sonnet 5 (Medium)
\begin{aiexample}[All partial derivatives exist, yet $f$ is not even continuous]
\label{ex:partial_derivatives_not_continuous}
% Generator: Claude Sonnet 5 (Medium)
Let $f : \mathbb{R}^2 \to \mathbb{R}$ be defined by
\[f(x,y) := \begin{cases} \dfrac{xy}{x^2+y^2} & (x,y) \neq (0,0), \\ 0 & (x,y) = (0,0). \end{cases}\]
Along each axis, $f$ vanishes identically ($f(x,0) = f(0,y) = 0$), so both partial derivatives at
the origin exist and equal $0$:
\[\partial_x f(0,0) = \lim_{s\to 0}\frac{f(s,0)-f(0,0)}{s} = 0, \qquad \partial_y f(0,0) = 0.\]
Yet $f$ is not even continuous at $(0,0)$: along the diagonal $y=x$, $f(x,x) = \tfrac{x^2}{2x^2} =
\tfrac12$ for all $x \neq 0$, so $f(x,x) \to \tfrac12 \neq 0 = f(0,0)$ as $x \to 0$. By
\Cref{prop:differentiable_implies_continuous}, $f$ is therefore not differentiable at $(0,0)$,
even though both partial derivatives exist there.
\end{aiexample}

% Supplement: Sascha Brack/Class Notes/Week_03_Notes_Monday.pdf, p. 7
\begin{remark}[The same expression, squared, behaves completely differently]
\label{rem:continuity_contrast_pair}
Sascha Brack pairs the function above with one that looks almost identical, as a test of whether
you can tell them apart:
\[f_1(x,y) := \frac{xy}{x^2+y^2}, \qquad\qquad f_2(x,y) := \frac{x^2y^2}{x^2+y^2},\]
both extended by $0$ at the origin. Substituting $x = r\cos\theta$, $y = r\sin\theta$ settles
both at once:
\[f_1 = \frac{r^2\cos\theta\sin\theta}{r^2} = \cos\theta\sin\theta,
\qquad\qquad
f_2 = \frac{r^4\cos^2\theta\sin^2\theta}{r^2} = r^2\cos^2\theta\sin^2\theta.\]
In $f_1$ every $r$ cancels: the function is constant along each ray and takes a different
constant on different rays, so no limit exists and $f_1$ is discontinuous at the origin. In
$f_2$ a factor $r^2$ survives, and since $|\cos^2\theta\sin^2\theta| \leq 1$ we get
$|f_2| \leq r^2 \to 0$ \emph{uniformly in $\theta$}, so $f_2$ is continuous at the origin.

The pair is worth doing once because it shows what the polar substitution is actually testing.
The question is never \qt{does the expression go to zero along each fixed ray?}, which $f_1$ does
not even manage. The question is \textbf{does the $\theta$-dependence get killed by a positive
power of $r$?} If the $r$'s cancel completely, the direction survives into the limit and there is
none.
This is the same test used in \cref{q:quiz_q4} on the repetition quiz, and the reason
step 3 of \cref{rem:differentiability_recipe} insists on \qt{uniformly in $\theta$}.
\end{remark}

% Sonnet 5 (Medium)
\begin{center}
\begin{tikzpicture}[>=Latex, scale=1.0]
  \draw[->] (-1.6,0) -- (1.6,0) node[right, font=\small] {$x$};
  \draw[->] (0,-1.6) -- (0,1.6) node[above, font=\small] {$y$};
  \draw[orange!90!black, thick, ->] (0,0) -- (1.3,0) node[below right, font=\scriptsize] {$f\equiv0$};
  \draw[orange!90!black, thick, ->] (0,0) -- (0,1.3) node[above left, font=\scriptsize] {$f\equiv0$};
  \draw[purple, thick, ->] (0,0) -- (1.1,1.1) node[above right, font=\scriptsize] {$f\equiv\tfrac12$};
  % Generator: Claude Opus 5 (High) --- caption added; the picture previously stood with no
  % text of its own, several paragraphs away from the example it illustrates.
  \node[below, font=\footnotesize, align=center] at (0,-1.75)
    {$f = \tfrac{xy}{x^2+y^2}$ near the origin: constant along every ray,\\
     but a \emph{different} constant on the diagonal than on the axes};
\end{tikzpicture}
\end{center}

% Transition: Claude Opus 5 (High) --- the diagram below arrived with no text at all around
% it, between a figure and a \subsection heading.
Sascha Brack's notes collect these relationships into one picture. Read it as a hierarchy: the
strongest condition is at the bottom, and every arrow points at something weaker.

% Supplement: Sascha Brack/Class Notes/Week_04_Notes_Monday_Updated.pdf, p. 4
\begin{center}
\begin{tikzcd}[row sep=3.5em, column sep=1.5em]
  & \begin{tabular}{c}\textbf{Differentiable}\\($Df_{x_0}$ exists)\end{tabular} \arrow[dl, Rightarrow] \arrow[dr, Rightarrow] & \\
  \begin{tabular}{c}\textbf{Partial derivatives}\\exist\end{tabular} & & \begin{tabular}{c}\textbf{Directional derivatives}\\exist\end{tabular} \\
  \begin{tabular}{c}\textbf{Partial derivatives}\\exist\\and are continuous\end{tabular} \arrow[u, Rightarrow] \arrow[uur, Rightarrow] \arrow[rr, Leftrightarrow] & & \begin{tabular}{c}\textbf{Directional derivatives}\\exist\\and are continuous\end{tabular} \arrow[u, Rightarrow] \arrow[uul, Rightarrow]
\end{tikzcd}
\end{center}

% Generator: Claude Opus 5 (High)
\begin{importantremark}[Not one of these arrows reverses]
The diagram is easy to misread as a set of equivalences, and it is nothing of the kind: every
$\Rightarrow$ in it is strict. \Cref{ex:partial_derivatives_not_continuous} already shows the
bottom-left arrow failing backwards, since there both partial derivatives exist at the origin
while $f$ is not even continuous. The other two failures, with their counterexamples, are collected in
\cref{sec:sufficient_condition_differentiability}, which redraws this same landscape as a single
vertical chain and attaches a witness to each broken link.

The one genuine equivalence is the horizontal arrow along the bottom, and it is worth seeing
why. Left to right: continuous partials make $f$ continuously differentiable
(\cref{thm:continuous_partials_imply_differentiable}), so every directional derivative exists
and equals $Df_x(v)$, which depends continuously on $x$. Right to left is immediate: take
$v = e_i$.
\end{importantremark}

\subsection{Examples}

% Transition: Claude Opus 5 (High)
The three examples that follow are the three shapes the Jacobian takes, and it is worth naming
them, because the layout of the matrix changes completely between them while the underlying
object, the linear map $DF_{x_0}$, does not.
\begin{itemize}
  \item \textbf{A curve} ($n=1$, $m$ arbitrary): one column, so $\Jac\gamma(t)$ is a
    \emph{column vector}, the velocity.
  \item \textbf{A scalar field} ($n$ arbitrary, $m=1$): one row, so $\Jac f(x)$ is a
    \emph{row vector}, namely the transposed gradient $(\nabla f)\transp$.
  \item \textbf{A vector field} (both $>1$): the full $m \times n$ matrix, whose columns are the
    partial derivatives and whose rows are the gradients of the components.
\end{itemize}
That the Jacobian of a scalar field is a row while the Jacobian of a curve is a column is not a
convention to memorise: it follows from \qt{$m$ rows, $n$ columns} and nothing else. Note the
transpose in the second bullet, which is the one place a real convention enters: the
\emph{gradient} $\nabla f$ is a column vector, so it is $\Jac f = (\nabla f)\transp$ that is the
row. Confusing the two is the commonest source of shape errors when the chain rule is applied
later in this chapter.

% Source: Corsin Nick/Class Notes/Week 4.pdf, p. 4
\begin{example}[Derivative of a curve]
\[\gamma : \mathbb{R}\to\mathbb{R}^2, \qquad t \mapsto \begin{pmatrix}\cos t \\ \sin t\end{pmatrix}\]
What is $\gamma'(\tfrac{\pi}{4}) := D\gamma_{\pi/4}$?
\[\gamma'\!\left(\tfrac{\pi}{4}\right) = \begin{pmatrix}\partial_t\gamma_1(\pi/4) \\ \partial_t\gamma_2(\pi/4)\end{pmatrix} = \begin{pmatrix}-1/\sqrt{2} \\ 1/\sqrt{2}\end{pmatrix}\]
\end{example}

\begin{aiexample}[Jacobian Matrix of Polar Coordinates]
% Generator: Gemini 3.6 Flash (Medium)
Let $f \colon (0,\infty) \times (0,2\pi) \to \mathbb{R}^2$ be the polar coordinate map $f(r,\theta) := (r\cos\theta, r\sin\theta)\transp$. The Jacobian matrix $\Jac f(r,\theta)$ is:
\[
\Jac f(r,\theta) = \begin{pmatrix} \dfrac{\partial f_1}{\partial r} & \dfrac{\partial f_1}{\partial \theta} \\[1ex] \dfrac{\partial f_2}{\partial r} & \dfrac{\partial f_2}{\partial \theta} \end{pmatrix} = \begin{pmatrix} \cos\theta & -r\sin\theta \\ \sin\theta & r\cos\theta \end{pmatrix}.
\]
Its determinant is $\det \Jac f(r,\theta) = r\cos^2\theta + r\sin^2\theta = r > 0$.
\end{aiexample}

\begin{center}
\begin{tikzpicture}[>=Latex, scale=1.3]
  \draw[->] (-1.4,0) -- (1.4,0) node[right] {$x$};
  \draw[->] (0,-1.4) -- (0,1.4) node[above] {$y$};
  \draw[thick] (0,0) circle (1.0);
  \coordinate (P) at (0.707, 0.707);
  \draw[dotted, orange!90!black, thick] (0,0) -- (P);
  \fill[orange!90!black] (P) circle (2pt) node[above right, font=\footnotesize] {$\gamma(\pi/4)$};
  \draw[orange!90!black, ultra thick, ->] (P) -- ++(-0.6, 0.6) node[above left, font=\small] {$\gamma'(\pi/4)$};
\end{tikzpicture}
\end{center}

This vector is the \textbf{velocity of the curve} \germanfor{Geschwindigkeit} at
$t = \tfrac{\pi}{4}$; its length $|\gamma'(t)|$ is the \textbf{speed}. We get that
\[\gamma\!\left(\tfrac{\pi}{4}+h\right) = \frac{1}{\sqrt{2}}\begin{pmatrix}1\\1\end{pmatrix} + \frac{h}{\sqrt{2}}\begin{pmatrix}-1\\1\end{pmatrix} + o(|h|).\]

\begin{ainote}
Corsin calls $\gamma'(\tfrac{\pi}{4})$ the \qt{speed} of the curve. Strictly, $\gamma'(t)$ is the
\emph{velocity}, a vector carrying both a direction and a magnitude, whereas the speed is the
scalar $|\gamma'(t)|$, here $|(-1,1)|/\sqrt2 = 1$. Since the display below uses the vector, the
distinction matters. Corrected above.
\end{ainote}

% Source: Corsin Nick/Class Notes/Week 4.pdf, p. 4
From this example, we see that the \textbf{column} vectors of the Jacobian are the partial
derivatives:
\[\Jac F(x) = \left(\frac{\partial F}{\partial x_1}\ \Big|\ \cdots\ \Big|\ \frac{\partial F}{\partial x_n}\right)\]
and therefore, in the canonical basis, with $v := (v_1,\dots,v_n)\transp$:
\[\partial_v F(x) = DF_x(v) = \Jac F(x)\,v = \sum_{i=1}^{n} v_i\,\frac{\partial F}{\partial x_i}.\]

% Source: Corsin Nick/Class Notes/Week 4.pdf, p. 5
\begin{example}[Scalar field gradient]
\[f : \mathbb{R}^3\to\mathbb{R}, \qquad (x,y,z) \mapsto 2x^2 + y^2 + 3z^2 - 2xyz\]
What is $\Jac f(x,y,z)$?
\[\Jac f(x,y,z) = (4x - 2yz,\ 2y - 2xz,\ 6z - 2xy) = \left(\frac{\partial f}{\partial x}, \frac{\partial f}{\partial y}, \frac{\partial f}{\partial z}\right) = (\nabla f)\transp\]
\end{example}

% Supplement: Sascha Brack/Class Notes/Week_04_Notes_Friday_Updated.pdf, p. 3
\begin{example}[Jacobian matrix of a vector field]
\label{ex:jacobian_vector_field}
Consider $f : \mathbb{R}^2 \to \mathbb{R}^3$ defined by
\[f(x,y) := \begin{pmatrix} x^2+y^2 \\ 2xy \\ e^x\sin(y) \end{pmatrix}.\]
The partial derivatives with respect to $x$ and $y$ are the column vectors:
\[\partial_1 f(x,y) = \begin{pmatrix} 2x \\ 2y \\ e^x\sin(y) \end{pmatrix}, \qquad \partial_2 f(x,y) = \begin{pmatrix} 2y \\ 2x \\ e^x\cos(y) \end{pmatrix}.\]
Because these partial derivatives exist everywhere and are continuous, $f$ is differentiable on all of $\mathbb{R}^2$. The differential $Df_{(x_0,y_0)}$ is represented by the $3 \times 2$ Jacobian matrix:
\[\Jac f(x_0,y_0) = \begin{pmatrix} 2x_0 & 2y_0 \\ 2y_0 & 2x_0 \\ e^{x_0}\sin(y_0) & e^{x_0}\cos(y_0) \end{pmatrix}.\]
\end{example}

From this, we see that the \textbf{rows} of the Jacobian are the \newterm{gradients}
\germanfor{Gradient}:
\[\Jac F(x) = \begin{pmatrix} \nabla F_1 \\ \hline \vdots \\ \hline \nabla F_m \end{pmatrix}\]

\begin{ainote}
Corsin's sentence reads ``the columns of the Jacobian are the gradients'', but the displayed
matrix stacks $\nabla F_1, \dots, \nabla F_m$ as rows. The display is the correct one, and it is
consistent with $(\nabla f)\transp$ just above. Corrected to ``rows''.
\end{ainote}

% Supplement: Sascha Brack/Class Notes/Week_04_Notes_Monday_Updated.pdf, pp. 10--11
\begin{example}[A function that is differentiable but not $C^1$]
\label{ex:differentiable_not_c1}
Let $f : \mathbb{R}^2 \to \mathbb{R}$ be defined by
\[
f(x,y) := \begin{cases}
  (x^2+y^2)\sin\left(\dfrac{1}{\sqrt{x^2+y^2}}\right), & \text{if } (x,y) \neq (0,0), \\
  0, & \text{if } (x,y) = (0,0).
\end{cases}
\]
We claim that $f$ is differentiable at $(0,0)$ with $Df_{(0,0)} = (0, 0)$, but the partial derivatives are not continuous at $(0,0)$. 

\textbf{Step 1: Differentiability at $(0,0)$.} We use the definition of the differential. Let $L = (0, 0)$. We need to show that
\[
  \lim_{h \to (0,0)} \frac{f(h_1, h_2) - f(0,0) - L\begin{pmatrix}h_1\\h_2\end{pmatrix}}{|h|} = 0.
\]
Plugging in the definitions and changing to polar coordinates $(h_1, h_2) = (r\cos\theta, r\sin\theta)$ with $r = |h| = \sqrt{h_1^2+h_2^2}$:
\[
  \lim_{r \to 0} \frac{r^2 \sin(1/r) - 0 - 0}{r} = \lim_{r \to 0} r \sin(1/r) = 0.
\]
The sine term is bounded between $-1$ and $1$, so the limit evaluates to $0$. Thus, $f$ is differentiable at $(0,0)$ and $Df_{(0,0)} = (0, 0)$, which implies $\partial_1 f(0,0) = 0$ and $\partial_2 f(0,0) = 0$.

\textbf{Step 2: Discontinuity of partial derivatives.} Let us compute $\partial_1 f(x,y)$ for $(x,y) \neq (0,0)$. Using the product and chain rules:
\[
  \partial_1 f(x,y) = 2x \sin\left(\frac{1}{\sqrt{x^2+y^2}}\right) - \frac{x}{\sqrt{x^2+y^2}} \cos\left(\frac{1}{\sqrt{x^2+y^2}}\right).
\]
Let us choose a sequence $(x_n, y_n) := (\frac{1}{n}, 0)$ which converges to $(0,0)$ as $n \to \infty$. Then
\[
  \lim_{n \to \infty} \partial_1 f\left(\frac{1}{n}, 0\right) = \lim_{n \to \infty} \left[ \frac{2}{n} \sin(n) - \cos(n) \right].
\]
The first term vanishes, but $\lim_{n \to \infty} \cos(n)$ does not exist. Thus, $\partial_1 f$ is not continuous at $(0,0)$, and so $f$ is not continuously differentiable ($C^1$) at $(0,0)$. (By symmetry, the same holds for $\partial_2 f$).
\end{example}

% Sonnet 5 (Medium)


% --- exercises relocated here from the dissolved sheet 4 ---

% Originally: content/week-04.tex, end of the chapter

\begin{aiexercise}[Differential and Linearization at $(1,0)$]
\label{ex:ai_differential_linearization}
% Generator: Gemini 3.6 Flash (Medium)
Let $f(x,y) := x^2 y + \sin(xy)$ defined on $\mathbb{R}^2$.
\begin{enumerate}[label=\textbf{(\alph*)}]
  \item Compute the gradient $\nabla f(x,y)$ and the total differential $Df(1,0)$.
  \item Find the linear approximation $L(x,y)$ of $f$ around the point $(1,0)$.
\end{enumerate}
\end{aiexercise}

% Quelle: Linus Lüchinger/Slides/Slides-03-19.pdf, slides 9--11
% Generator: Claude Opus 5 (High)
\subsection{Differentiating with respect to a matrix}
\label{sec:differentiating_wrt_matrix}

Nothing in \cref{def:differentiable} requires the domain to be $\mathbb{R}^n$ with $n$ small, or
even to be written as a column of numbers. A function of an $n\times m$ matrix is a function of
$nm$ real variables, and its partial derivatives $\partial f/\partial A_{ij}$ are defined exactly
as before, by varying one entry and freezing the rest.

% Reclassified from ainote to remark: motivation for the mathematics, not commentary on the
% transcription. Linus spends part of a session making this argument.
\begin{remark}[Why this is worth a subsection]
Taking derivatives with respect to a matrix argument can look like a formal exercise with no
purpose. It is not, and the reason is worth recording: \textbf{it is how neural networks are
trained.} The parameters of such a network are matrices. The function being minimised sends those
matrices to a single number measuring how badly the network fits the data, and training consists
of computing the gradient with respect to those matrices and then stepping a little way against
it, because $-\nabla f$ points in the direction of steepest decrease. Every step of that loop is
the material of this chapter.
\end{remark}

The simplest instance is already worth doing by hand, and it is exactly linear regression.

% Generator: Claude Opus 5 (High)
\begin{aiexample}[The gradient of a least-squares fit]
\label{ex:least_squares_matrix_gradient}
Fix data matrices $X \in \mathbb{R}^{m\times k}$ (the $k$ inputs, each a column in
$\mathbb{R}^m$) and $Y \in \mathbb{R}^{n\times k}$ (the $k$ desired outputs). The model applies
an unknown matrix $A \in \mathbb{R}^{n\times m}$ and adds an unknown offset
$b \in \mathbb{R}^n$ to every column, so its total error is
\[f(A,b) := \big\lVert \underbrace{AX + b\mathbf{1}\transp - Y}_{=:\,R} \big\rVert^2,
  \qquad \mathbf{1} := (1,\dots,1)\transp \in \mathbb{R}^k,\]
where $\lVert\cdot\rVert$ is the Hilbert--Schmidt norm of \cref{sec:structured_spaces}, so that
$f(A,b) = \sum_{p,q} R_{pq}^2$ is the total squared error over all $k$ data points. Writing out
one entry of the residual,
\[R_{pq} = \sum_{s=1}^{m} A_{ps}X_{sq} + b_p - Y_{pq},\]
so that $\partial R_{pq}/\partial A_{ij} = \delta_{pi}X_{jq}$ and
$\partial R_{pq}/\partial b_i = \delta_{pi}$. The chain rule then gives
\begin{align}
  \frac{\partial f}{\partial A_{ij}} &= 2\sum_{p,q} R_{pq}\,\delta_{pi}X_{jq}
    = 2\sum_{q} R_{iq}X_{jq} = 2\big(RX\transp\big)_{ij}, \nonumber \\
  \frac{\partial f}{\partial b_i} &= 2\sum_{p,q} R_{pq}\,\delta_{pi}
    = 2\sum_{q} R_{iq} = 2\big(R\mathbf{1}\big)_i. \nonumber
\end{align}
Collecting the entries back into arrays of the same shape as the variables,
\[\nabla_{\!A} f = 2\big(AX + b\mathbf{1}\transp - Y\big)X\transp, \qquad
  \nabla_{\!b} f = 2\big(AX + b\mathbf{1}\transp - Y\big)\mathbf{1}.\]
The gradient with respect to $b$ is twice the vector of row sums of the residual, and the
gradient with respect to $A$ is the residual correlated against the inputs. Both are entirely
explicit, and both were computed with nothing beyond partial differentiation and the chain rule.
\end{aiexample}

% Reclassified from ainote to remark: a caution about how the identity is read is mathematics.
\begin{remark}[A notational caution]
One point trips people up here. Since $f$ above has domain
$\mathbb{R}^{n\times m}\times\mathbb{R}^n$ of dimension $nm+n$, its Jacobian is formally a
$1\times(nm+n)$ row, and $\partial_v f(A,b) = \Jac f(A,b)\,v$ would require flattening the matrix
direction $v$ into a column of length $nm$. That flattening is pure bookkeeping and is almost
never performed: one keeps $\nabla_{\!A}f$ shaped like $A$, as above. The identity
$\partial_vf = \Jac f\cdot v$ of \cref{sec:jacobi_matrix} still holds, and is simply being read
in a reshaped coordinate system.
\end{remark}

% Originally: content/week-04.tex, \exercisesheet{4}
% Source: exercises/Ex4_Analysis2_eng.pdf

\begin{exercise}[Derivative computation \textnormal{(important)}]
\label{ex:4.1}
Consider the function $u : (x,y) \mapsto x^{\sin(y)}$, defined for
$(x,y) \in (0,\infty)\times\mathbb{R} \subset \mathbb{R}^2$. Compute $\partial_x u$ and
$\partial_y u$.
\exinfo{This exercise is Problem 4.1 of Problem Sheet 4, Analysis II, Spring Semester 2026.}
\end{exercise}

% Source: old_exams/examHS24.pdf (13 February 2025), Wahr/Falsch-Fragen, Aufgabe 4, p. 3
% Extractor: Gemini 3.1 Pro (High)
\begin{exercise}[True or False \textnormal{(piecewise differentiability)}]
\label{ex:examHS24_TF_diff_piecewise}
Consider the function $f \colon \mathbb{R}^2 \to \mathbb{R}$ given by
\[
f(x,y) = \begin{cases}
y^2 \cos x & \text{if } y \ge 0, \\
0 & \text{if } y < 0.
\end{cases}
\]
Decide whether each of the following statements is true or false.
\begin{enumerate}[label=\textbf{(\alph*)}]
  \item $f \in C^0(\mathbb{R}^2)$.
  \item $\partial_y f \in C^0(\mathbb{R}^2)$.
  \item $\partial_y f \in C^1(\mathbb{R}^2)$.
  \item $f \in C^2(\mathbb{R}^2)$.
  \item $\frac{\partial^2 f}{\partial x^2}(0,0)$ exists.
\end{enumerate}
\exinfo{This exercise is taken from the exam of 13 February 2025 (Prof.\ Joaquim Serra), where it appears as the fourth
block of true/false questions.}
\end{exercise}
